\documentclass{beamer}
\usetheme{Madrid}
\usecolortheme{seahorse}

\title{Analisi Prezzi Aeromobili}
\subtitle{Modelli Lineari per la Predizione del Prezzo}
\author{Statistical Learning Project}
\date{\today}

\usepackage[utf8]{inputenc}
\usepackage[italian]{babel}
\usepackage{graphicx}
\usepackage{amsmath}
\usepackage{booktabs}

\begin{document}

\frame{\titlepage}

% 1. Introduzione a dataset
\begin{frame}{Introduzione al Dataset}
\begin{itemize}
    \item \textbf{Dataset}: Prezzi di aeromobili
    \item \textbf{Dimensioni}: 507 osservazioni, 16 variabili
    \item \textbf{Target}: Price (USD)
    \item \textbf{Fonte}: Dati reali di mercato aeronautico
\end{itemize}

\vspace{1cm}
\begin{center}
\textbf{Obiettivo}: Predire il prezzo degli aeromobili utilizzando caratteristiche tecniche
\end{center}
\end{frame}

% 2. Features e obiettivi
\begin{frame}{Features e Obiettivi}
\textbf{Variabili Principali:}
\begin{columns}
\begin{column}{0.5\textwidth}
\begin{itemize}
    \item cruise\_speed
    \item max\_speed  
    \item fuel\_tank
    \item engine\_power
    \item wing\_span
    \item landing\_distance
    \item empty\_weight
    \item range
\end{itemize}
\end{column}
\begin{column}{0.5\textwidth}
\begin{itemize}
    \item manufacturer
    \item engine\_type
    \item aircraft\_type
    \item cockpit\_type
    \item fuel\_type
    \item category
    \item seating\_capacity
    \item \textcolor{red}{\textbf{price}} (target)
\end{itemize}
\end{column}
\end{columns}

\vspace{0.5cm}
\textbf{Obiettivo}: Costruire un modello lineare multiplo per predire il prezzo
\end{frame}

% 3. Grafici plot rispetto a price
\begin{frame}{Relazioni con il Prezzo}
\begin{center}
\textbf{Scatter Plots delle Variabili Numeriche vs Price}
\end{center}

\begin{itemize}
    \item \textbf{Correlazioni positive forti}:
    \begin{itemize}
        \item cruise\_speed vs price
        \item max\_speed vs price
        \item fuel\_tank vs price
        \item engine\_power vs price
    \end{itemize}
    \item \textbf{Relazioni non lineari evidenti}
    \item \textbf{Presenza di outliers} nei dati
\end{itemize}
\end{frame}

% 4. Plot distribuzioni
\begin{frame}{Distribuzioni delle Variabili}
\textbf{Analisi delle Distribuzioni:}

\begin{itemize}
    \item \textbf{Price}: Fortemente asimmetrica a destra
    \item \textbf{Fuel\_tank}: Distribuzione skewed
    \item \textbf{Engine\_power}: Asimmetria positiva
    \item \textbf{Landing\_distance}: Skewness elevata
    \item \textbf{Empty\_weight}: Coda lunga a destra
    \item \textbf{Range}: Distribuzione non normale
    \item \textbf{Wing\_span}: Asimmetria moderata
\end{itemize}

\vspace{0.5cm}
\textcolor{red}{\textbf{Problema}}: Le distribuzioni non normali violano le assunzioni del modello lineare
\end{frame}

% 5. Heatmap
\begin{frame}{Heatmap Correlazioni}
\begin{center}
\textbf{Matrice di Correlazione - Dati Originali}
\end{center}

\textbf{Correlazioni Elevate:}
\begin{itemize}
    \item price - cruise\_speed: 0.85
    \item price - max\_speed: 0.82
    \item price - fuel\_tank: 0.89
    \item price - engine\_power: 0.91
    \item landing\_distance - engine\_power: 0.95 (\textcolor{red}{multicollinearità})
\end{itemize}

\vspace{0.5cm}
\textcolor{orange}{\textbf{Evidenza di multicollinearità}} tra alcune variabili
\end{frame}

% 6. Log Trasformazione e plot distribuzioni
\begin{frame}{Trasformazione Logaritmica}
\textbf{Variabili Trasformate con Log:}
\begin{itemize}
    \item price → log(price)
    \item fuel\_tank → log(fuel\_tank)
    \item engine\_power → log(engine\_power)
    \item landing\_distance → log(landing\_distance)
    \item empty\_weight → log(empty\_weight)
    \item range → log(range)
    \item wing\_span → log(wing\_span)
\end{itemize}

\vspace{0.5cm}
\textbf{Risultato}: Distribuzioni più simmetriche e normali
\end{frame}

% 7. Heatmap post-log
\begin{frame}{Heatmap Post-Trasformazione}
\begin{center}
\textbf{Matrice di Correlazione - Dati Log-Trasformati}
\end{center}

\textbf{Miglioramenti:}
\begin{itemize}
    \item Correlazioni più lineari
    \item Riduzione della multicollinearità
    \item Relazioni più stabili
    \item Distribuzioni normalizzate
\end{itemize}

\vspace{0.5cm}
\textcolor{green}{\textbf{Beneficio}}: Assunzioni del modello lineare meglio soddisfatte
\end{frame}

% 8. Modello multiplo lineare
\begin{frame}{Modello Lineare Multiplo}
\textbf{Performance del Modello:}
\begin{itemize}
    \item \textbf{R²}: 0.865 (86.5\% varianza spiegata)
    \item \textbf{RMSE}: \$382,447
    \item \textbf{Variabili significative}: 10/13 (p < 0.05)
\end{itemize}

\textbf{Cross-Validation:}
\begin{itemize}
    \item Train/Test Split (80-20): R² = 0.849
    \item 10-Fold CV: R² = 0.84-0.86
\end{itemize}

\textbf{Variabili Non Significative:}
\begin{itemize}
    \item aircraft\_type\_Light
    \item engine\_type\_Turbo-prop  
    \item manufacturer\_Cessna
    \item manufacturer\_Piper
\end{itemize}
\end{frame}

\begin{frame}{Conclusioni}
\begin{center}
\textbf{Risultati Principali}
\end{center}

\begin{itemize}
    \item \textcolor{green}{\textbf{Modello robusto}} con R² = 86.5\%
    \item \textcolor{green}{\textbf{Cross-validation}} conferma stabilità
    \item \textcolor{blue}{\textbf{Trasformazione logaritmica}} essenziale
    \item \textcolor{orange}{\textbf{Multicollinearità}} ancora presente
\end{itemize}

\vspace{1cm}
\begin{center}
\textbf{Il modello lineare multiplo è efficace per la predizione dei prezzi degli aeromobili}
\end{center}
\end{frame}

\end{document}
